\section{Penutup}

% Slide 31: Bab 5.1 - Kesimpulan (1)
\begin{frame}{Bab 5.1 - Kesimpulan: Formulasi EPG-Hamiltonian}
    \begin{itemize}
        \item \textbf{Integrasi Teori-Fisis:} Formulasi Markowitz ke dalam kerangka \textit{Exact Potential Game} (EPG) berhasil memetakan utilitas ekonomi menjadi sistem Hamiltonian fisis secara konsisten.
        \item \textbf{Presisi Parameter:}
        \begin{itemize}
            \item \textbf{Log Return:} Menjaga simetri interaksi risiko ($J_{ij}$) tanpa distorsi matematis.
            \item \textbf{Simple Return:} Mempertahankan linearitas bias ekspektasi imbal hasil ($h_i$).
        \end{itemize}
        \item \textbf{Integritas Solusi:} Implementasi suku penalti anggaran ($A$) efektif mencegah sirkuit kuantum terjebak pada status yang melanggar batasan investasi.
    \end{itemize}
\end{frame}

% Slide 32: Bab 5.1 - Kesimpulan (2)
\begin{frame}{Bab 5.1 - Kesimpulan: Efektivitas Warm-Start Nash}
    \begin{itemize}
        \item \textbf{Akselerasi Konvergensi:} Penerapan ekuilibrium nash sebagai inisialisasi \textit{warm-start} meningkatkan stabilitas konvergensi algoritma VQE secara signifikan.
        \item \textbf{Mitigasi Barren Plateaus:} Injeksi solusi strategis memandu sirkuit melewati area kelenyapan varians gradien.
        \item \textbf{Efisiensi Sumber Daya:} Algoritma mampu mencapai \textit{ground state} pada kedalaman sirkuit dangkal ($L=2$ hingga $L=3$), mereduksi beban komputasi sirkuit kuantum.
    \end{itemize}
\end{frame}

% Slide 33: Bab 5.1 - Kesimpulan (3)
\begin{frame}{Bab 5.1 - Kesimpulan: Kapabilitas dan Performa}
    \begin{itemize}
        \item \textbf{Ketangguhan Stokastik:} Sinergi arsitektur HEA dan optimizer SPSA menunjukkan kapabilitas tinggi dalam menangani volatilitas pasar modal.
        \item \textbf{Agilitas Rebalancing:} Estimasi gradien simultan memungkinkan respons cepat terhadap pergeseran ekuilibrium pasar.
        \item \textbf{Superioritas Empiris:} Strategi GT-VQE menghasilkan pertumbuhan modal paling unggul dengan \textbf{Sharpe Ratio 1,0107} (N=4), melampaui performa algoritma klasik SLSQP.
    \end{itemize}
\end{frame}

% Slide 34: Bab 5.2 - Saran Pengembangan
\begin{frame}{Bab 5.2 - Saran Pengembangan}
    \begin{enumerate}
        \item \textbf{Parameter Noise:} Integrasi model dekoherensi qubit dan \textit{gate fidelity} untuk implementasi pada perangkat keras kuantum riil.
        \item \textbf{Eksplorasi Arsitektur:} Pengujian variasi sirkuit \textit{ansatz} lain dan algoritma optimasi hibrida untuk meningkatkan laju konvergensi.
        \item \textbf{Analisis Krisis:} Perluasan data pada periode fenomena \textit{Black Swan} guna menguji reliabilitas algoritma dalam mempertahankan stabilitas ekuilibrium di tengah guncangan ekstrem.
    \end{enumerate}
    \vspace{0.8cm}
    \begin{center}
        \huge \textbf{Terima Kasih} \\
        \medskip
        \large Ada Pertanyaan?
    \end{center}
\end{frame}
